\documentclass[../DoAn.tex]{subfiles}
\begin{document}

\section{Thiết kế kiến trúc}
\subsection{Lựa chọn kiến trúc phần mềm}
\label{subsection:4.1}
Phần mềm được xây dựng dựa trên kiến trúc client–server. Trong kiến trúc này, phía client chịu trách nhiệm tương tác với người dùng, gửi yêu cầu và hiển thị kết quả, còn phía server đảm nhiệm xử lý nghiệp vụ, truy xuất dữ liệu và trả về phản hồi. Kiến trúc này giúp tách biệt rõ ràng giữa giao diện và logic xử lý, từ đó tăng khả năng bảo trì, mở rộng và phát triển hệ thống.

Hình \ref{fig:Fig0} mô tả kiến trúc tổng quan của ứng dụng, dựa trên kiến trúc client-server. Cụ thể gồm 2 phần frontend (client) với các thành phần giao diện, hiển thị và backend (server) xử lý các yêu cầu.
\begin{figure}[H]
    \centering
    \includegraphics[width=\textwidth]{Hinhve/Kiến trúc tổng quan.png}
    \caption{Kiến trúc tổng quan của hệ thống}
    \label{fig:Fig0}
\end{figure}

Phía frontend gồm giao diện chat được custom từ giao diện chat của thư viện Chainlit, cho phép đặt câu hỏi, xem câu trả lời và theo dõi Agent trace. Giao diện visualize đồ thị tri thức là trang web riêng dùng thư viện vis-network; người dùng mở qua link trong câu trả lời để xem phần đồ thị tri thức được retriever truy xuất. Frontend duy trì HTTP/WebSocket với server; trang viz gọi HTTP GET tới server để lấy dữ liệu vẽ đồ thị tri thức.

Phía backend có nhiệm vụ điều phối Agentic Workflow, quản lý session hỏi đáp, xử lý truy vấn của người dùng, gọi API LLM, truy vấn đồ thị tri thức và sinh câu trả lời có căn cứ pháp lý. Các thành phần chính: (i) Chat Orchestrator (Luồng điều phối chat) nhận câu hỏi từ người dùng, duy trì lịch sử hội thoại của đoạn chat, gọi Router Agent, xử lý và hợp nhất context nhận được từ các retriever, stream câu trả lời từ LLM (gpt-4.1); (ii) Router Agent gọi API LLM (gpt-4o) theo cơ chế tool-calling để chọn (các) retriever; (iii) Retriever Agents gồm 20 retrievers được chia theo dạng câu hỏi pháp lý, mỗi retriever phân loại template cần dùng để trả lời cho câu hỏi, trích xuất tham số bằng LLM (gpt-4o) rồi thực hiện truy vấn Cypher; (iv) Neo4j Knowledge Graph để lưu đồ thị tri thức Luật Hôn nhân và Gia đình 2014; (v) Session Storage bao gồm bộ nhớ Chainlit in-memory storage để lưu trữ các dữ liệu gắn liền với phiên chat hiện tại; PostgreSQL (Chainlit Data Layer) lưu trữ lịch sử hỏi đáp, input/output của agent, context truy xuất, câu trả lời cuối cùng và các thông tin cảnh báo nhằm khôi phục lại đoạn hội thoại cũ, kiểm thử, truy vết và đánh giá hệ thống.

\subsection{Thiết kế tổng quan}
\label{subsection:4.2}

% \subsubsection{Thiết kế gói phía server}
% \label{subsubsection:4.2.1}
\begin{figure}[H]
    \centering
    \includegraphics[width=0.72\textwidth]{Hinhve/Sơ đồ gói server.png}
    \caption{Biểu đồ phụ thuộc gói phía server}
    \label{fig:Fig1}
\end{figure}
Phía server được tổ chức thành bốn gói được thể hiện ở Hình \ref{fig:Fig1} bao gồm (i) Gói presentation để xây dựng các chức năng xử lý sự kiện hội thoại, xác thực OAuth và phiên đăng nhập guest, quản lý projects hội thoại. (ii) Gói application đảm nhiệm xử lý logic và điều phối Agentic Workflow. (iii) Gói adapter chịu trách nhiệm tương tác với cơ sở dữ liệu Neo4j để truy truy xuất đồ thị tri thức, quản lý lịch sử hội thoại trên cơ sở dữ liệu PostgreSQL, cấu hình API LLM, trực quan hóa đồ thị tri thức , nhận diện người dùng guest. (iv) Gói utils là gói dùng chung cho toàn hệ thống, cung cấp các hàm chuẩn hóa context pháp lý cho LLM, các hàm xử lý chuỗi và thời gian.

% \subsubsection{Thiết kế gói phía client}
% \label{subsubsection:4.2.2}
\begin{figure}[H]
    \centering
    \includegraphics[width=\textwidth]{Hinhve/Sơ đồ gói client.png}
    \caption{Biểu đồ phụ thuộc gói phía client}
    \label{fig:Fig1_client}
\end{figure}

Phía client được xây dựng dựa trên frontend của Chainlit phiên bản 2.11.1 có tùy chỉnh lại. Ngoài giao diện chatbot chính dựa trên frontend của Chainlit, hệ thống còn có một trang web tĩnh độc lập dùng thư viện vis-network để trực quan hóa đồ thị tri thức và người dùng có thể mở liên kết tới trang web tĩnh tĩnh này trong câu trả lời của chatbot. Hình \ref{fig:Fig1_client} mô tả biểu đồ gói tổng quan phía client cùng phân loại các gói được tùy chỉnh/giữ nguyên từ frontend Chainlit. Các gói bao gồm (i) Gói pages chứa các màn hình định tuyến chính. (ii) Gói components gom các thành phần UI tái sử dụng, ở đây tùy chỉnh lại LeftSidebar, thêm xử lý cho projects hỏi đáp (iii) Gói api cung cấp api và WebSocket giao tiếp với Chainlit server. (iv) Gói lib chứa các tiện ích chia sẻ giữa pages và components, gồm các formatter và auth.ts cung cấp hàm canManageConversations dùng để ẩn quản lý hội thoại đối với guest. (v) Gói types định nghĩa các kiểu TypeScript dùng chung trong toàn ứng dụng, gói này được mở rộng bằng types/projects.ts cho chức năng project. (vi) Các gói hạ tầng state, contexts, hooks và i18n được giữ nguyên từ Chainlit frontend, phục vụ quản lý trạng thái, context react, custom hooks và đa ngôn ngữ.

% % \subsection{Thiết kế chi tiết gói}
% % Sinh viên thiết kế và lần lượt vẽ biểu đồ thiết kế cho từng package, hoặc một nhóm các package liên quan để giải quyết một vấn đề gì đó. Khi vẽ thiết kế gói, sinh viên chỉ cần đưa tên lớp, không cần chỉ ra các thành viên phương thức và thuộc tính. SV tham khảo ví dụ minh họa trong Hình \ref{fig:Fig2}.

% % Sinh viên cần vẽ rõ ràng quan hệ giữa các lớp trong biểu đồ. Các quan hệ bao gồm: phụ thuộc (dependency), kết hợp (association), kết tập (aggregation), hợp thành (composition), kế thừa (inheritance), và thực thi (implementation). Các quan hệ này đều đã được minh họa trong \ref{fig:Fig2}.

% % Sau khi vẽ hình minh họa, sinh viên cần giải thích ngắn gọn về thiết kế của mình. 

% \begin{figure}[H]
%     \centering
%     \includegraphics[width=0.72\textwidth]{Hinhve/Picture2.png}
%     \caption{Ví dụ thiết kế gói}
%     \label{fig:Fig2}
% \end{figure}

\section{Thiết kế chi tiết}
\subsection{Thiết kế giao diện}

Giao diện ứng dụng được xây dựng nhờ thư viện chainlit, hỗ trợ responsive trên các màn hình thiết bị khác nhau. Ngoài ra ứng dụng cũng xây dựng giao diện hiển thị đồ thị tri thức. Màu sắc được sử dụng trong giao diện chatbot và trang visualize đồ thị tri thức được thể hiện trong Bảng \ref{table:ui_color_palette}.

\begin{table}[H]
\centering
\begin{tabularx}{\textwidth}{|>{\raggedright\arraybackslash}p{3.5cm}|>{\raggedright\arraybackslash}p{3.0cm}|X|}
\hline
\textbf{Thành phần} & \textbf{Mã màu (Hex)} & \textbf{Mô tả} \\ \hline
Nền vùng hội thoại & \#212121 & Nền chính của màn hình trao đổi. \\ \hline
Thanh điều hướng (sidebar) & \#171717 & Nền tối hơn của khu vực quản lý dự án, lịch sử hội thoại và các chức năng điều hướng. \\ \hline
Thẻ và hộp nhập liệu & \#2E2E2E & Nền của thẻ nội dung, hộp nhập câu hỏi và các vùng nổi trên nền hội thoại chính. \\ \hline
Chữ chatbot và trang visualize đồ thị & \#EDEDED, \#B5B5B5 & Màu chữ chính và phụ dùng cho mô tả, gợi ý nhập liệu và thông tin ít quan trọng hơn. \\ \hline
Đường viền & \#434242 & Phân tách các hộp nhập, thẻ, đường biên. \\ \hline
Các nút hành động chính của chatbot & \#F5145F & Màu cho logo, ảnh đại diện, nút gửi và link đến trang visualize đồ thị tri thức. \\ \hline
Nền trang visualize đồ thị & \#121212, \#161616 & Lần lượt là nền tổng thể và nền vùng vẽ đồ thị. \\ \hline
Nút căn cứ pháp lý & \#4C8BF5, \#2E8B57, \#FFD700, \#FF0000 & Màu xanh lam cho căn cứ pháp lý chính, màu xanh lá cây cho căn cứ pháp lý bổ trợ, màu vàng cho căn cứ pháp lý thay thế hoặc bãi bỏ, màu đỏ cho căn cứ pháp lý mâu thuẫn. \\ \hline
Nút ngữ nghĩa & \#9013FE & Màu tím cho các nút ngữ nghĩa liên kết với căn cứ pháp lý. \\ \hline
Cạnh liên kết đồ thị & \#555555 & Màu xám cho đường nối và mũi tên quan hệ. \\ \hline
\end{tabularx}
\caption{Bảng màu chính của giao diện chatbot và trang visualize đồ thị tri thức}
\label{table:ui_color_palette}
\end{table}

Hình \ref{fig:user_flow_chatbot} mô tả user flow template của ứng dụng chatbot.

\begin{figure}[H]
    \centering
    \includegraphics[width=\textwidth]{Hinhve/User flow template chatbot.png}
    \caption{User flow template của ứng dụng chatbot}
    \label{fig:user_flow_chatbot}
\end{figure}

Màn hình đăng nhập được thiết kế đơn giản, với logo chính của ứng dụng chiếm nửa màn hình bên tay phải và khu vực xác thực ở vị trí bên tay trái và người dùng có thể đăng nhập thông qua các dịch vụ xác thực của bên thứ 3 mà Chainlit hỗ trợ (Ví dụ: Google OAuth). Thiết kế giao diện màn hình đăng nhập được minh họa trong Hình \ref{fig:ui_login}.

\begin{figure}[H]
    \centering
    \includegraphics[width=\textwidth]{Hinhve/Login mockup screen.png}
    \caption{Thiết kế giao diện màn hình đăng nhập}
    \label{fig:ui_login}
\end{figure}

Sau khi tiến hành xác thực thành công, người dùng được chuyển hướng đến màn hình giao tiếp chính của ứng dụng. Bố cục không gian được chia thành ba khu vực chính: thanh điều hướng (sidebar) và quản lý lịch sử hội thoại nằm bên trái, vùng hiển thị nội dung chat chiếm vị trí trung tâm, và thanh nhập liệu được cố định ở cạnh dưới màn hình. Các agent được gọi được hiển thị ngay phía trên vùng chứa tin nhắn trả lời, khi nhấn vào biểu tượng dropdown sẽ hiển thị chi tiết input, output của các agent. Vùng hiển thị các agent sẽ luôn tồn tại ngay cả khi câu trả lời cuối cùng của chatbot đã được đưa ra hoặc người dùng chuyển sang phiên hội thoại khác. Đường link đến trang visualize đồ thị tri tri thức và nút bấm xem cảnh báo hiệu lực hiển thị phía dưới cùng của câu trả lời. Màn hình trao đổi chính được trình bày trong Hình \ref{fig:ui_main_chat}.

\begin{figure}[H]
    \centering
    \includegraphics[width=\textwidth]{Hinhve/Giao diện màn hình trao đổi chính.png}
    \caption{Thiết kế giao diện màn hình trao đổi chính}
    \label{fig:ui_main_chat}
\end{figure}

Một trong những yêu cầu thiết kế cốt lõi của công cụ là khả năng hiển thị minh bạch luồng hoạt động của các hệ thống đa tác nhân (multi-agent) và đồ thị tri thức. Để giải quyết bài toán này, giao diện được tích hợp các khối thông tin có thể thu gọn ngay bên dưới mỗi câu trả lời của hệ thống. Thiết kế này cho phép người dùng quan sát chi tiết quá trình tác nhân phân tích câu hỏi, kiểm tra ngữ cảnh pháp lý được truy xuất từ đồ thị tri thức, và theo dõi dữ liệu đầu ra mà không làm ảnh hưởng đến trải nghiệm đọc câu trả lời cuối cùng. Giao diện hiển thị luồng xử lý của Agent được minh họa trong Hình \ref{fig:ui_agent_trace} và giao diện visualize đồ thị tri thức được minh họa trong Hình \ref{fig:ui_visualize_graph}.

\begin{figure}[H]
    \centering
    \includegraphics[width=\textwidth]{Hinhve/Màn hình hiển thị luồng xử lý agent và câu trả lời.png}
    \caption{Thiết kế giao diện hiển thị luồng xử lý và truy xuất dữ liệu của Agent}
    \label{fig:ui_agent_trace}
\end{figure}

\begin{figure}[H]
    \centering
    \includegraphics[width=\textwidth]{Hinhve/Màn hình hiển thị đồ thị tri thức.png}
    \caption{Thiết kế giao diện hiển thị đồ thị tri thức}
    \label{fig:ui_visualize_graph}
\end{figure}
% \subsection{Thiết kế lớp}
% Trong phần này, ĐATN sẽ trình bày thiết kế chi tiết một số module chủ đạo. 

% Đầu tiên là module main thuộc package Presentation. Module này đóng vai trò tiếp nhận câu hỏi từ người dùng và điều phối các agent. Hình \ref{fig:main_detail}, Bảng \ref{table:main_attributes}
% và Bảng \ref{table:main_functions} mô tả chi tiết các biến và các hàm thuộc module này.
% \begin{figure}[H]
%     \centering
%     \includegraphics[width=\textwidth]{Hinhve/thiết kế chi tiết cho module main.png}
%     \caption{Thiết kế chi tiết module main}
%     \label{fig:main_detail}
% \end{figure}

% \begin{table}[H]
% \centering
% \caption{Đặc tả các biến của module main}
% \begin{tabular}{|l|p{10cm}|}
% \hline
% \textbf{Biến} & \textbf{Ý nghĩa} \\ \hline
% tools & Lưu trữ danh sách các retriever của hệ thống, bao gồm tên, mô tả chức năng (description) và hàm thực thi tương ứng. \\ \hline
% main\_prompt & Chứa prompt định nghĩa các quy tắc lập luận cho LLM khi đưa ra câu trả lời  \\ \hline
% session\_history & Quản lý và lưu trữ dữ liệu hội thoại trong phiên làm việc hiện tại để cung cấp ngữ cảnh đầy đủ cho LLM khi đưa ra câu trả lời. \\ \hline
% \end{tabular}
% \label{table:main_attributes}
% \end{table}

% \begin{table}[H]
% \centering
% \caption{Đặc tả các hàm của module main}
% \begin{tabular}{|l|p{10cm}|}
% \hline
% \textbf{Hàm} & \textbf{Ý nghĩa} \\ \hline
% oauth\_callback() & Xử lý xác thực qua Google/Github và ánh xạ dữ liệu vào hệ thống lưu trữ PostgreSQL. \\ \hline
% on\_chat\_start() & Khởi tạo phiên làm việc, thiết lập lịch sử hội thoại mới và gửi thông điệp chào mừng. \\ \hline
% on\_chat\_resume() & Khôi phục trạng thái hội thoại từ cơ sở dữ liệu khi người dùng truy cập lại phiên cũ. \\ \hline
% main() & Phương thức điều phối thực hiện toàn bộ quy trình từ phân tích đến sinh câu trả lời. \\ \hline
% \end{tabular}
% \label{table:main_functions}
% \end{table}

% Tiếp theo sẽ trình bày về Module Router.

% Mô-đun router thuộc gói application đảm nhận vai trò của một Tác nhân định tuyến (Router Agent) thông minh. Chức năng cốt lõi của thành phần này là bóc tách các câu hỏi đa ý của người dùng và kích hoạt chính xác các retrievers cần thiết để lấy về ngữ cảnh pháp lý từ đồ thị tri thức nhằm trả lời cho câu hỏi của người dùng. Hình \ref{fig:router_detail}, Bảng  \ref{table:router_attributes} và bảng \ref{table:router_methods} mô tả chi tiết các biến và hàm trong module.

% \begin{figure}[H]
%     \centering
%     \includegraphics[width=\textwidth]{Hinhve/thiết kế chi tiết cho module router.png}
%     \caption{Thiết kế chi tiết module router}
%     \label{fig:router_detail}
% \end{figure}

% \begin{table}[H]
% \centering
% \begin{tabular}{|l|p{10cm}|}
% \hline
% \textbf{Biến} & \textbf{Ý nghĩa} \\ \hline
% tool\_picker\_prompt & Chứa hệ thống các quy tắc bắt buộc để LLM thực hiện như việc phân tích đa ý và lựa chọn công cụ, đảm bảo không gọi lặp lại một công cụ nhiều lần. \\ \hline
% \end{tabular}
% \caption{Đặc tả thuộc tính của mô-đun Router}
% \label{table:router_attributes}
% \end{table}

% \begin{table}[H]
% \centering
% \begin{tabular}{|l|p{10cm}|}
% \hline
% \textbf{Hàm} & \textbf{Ý nghĩa} \\ \hline
% handle\_tool\_calls() & Phương thức bất đồng bộ thực hiện việc gọi các hàm chức năng của công cụ, truyền tham số là câu truy vấn của người dùng và ghi lại kết quả truy vấn của từng công cụ nhờ hook của Chainlit\\ \hline
% tool\_choice() & Thực hiện việc liên kết (binding) danh sách công cụ với mô hình ngôn ngữ lớn và gọi mô hình để nhận về danh sách các công cụ cần kích hoạt. \\ \hline
% route\_question() & Hàm đầu vào chính của mô-đun, chịu trách nhiệm tổng hợp lịch sử hội thoại, câu hỏi đã chuẩn hóa để điều phối quy trình định tuyến và truy xuất dữ liệu. \\ \hline
% \end{tabular}
% \caption{Đặc tả các phương thức của mô-đun Router}
% \label{table:router_methods}
% \end{table}

% Đồ án thực hiện vẽ biểu đồ trình tự cho hai Use Case quan trọng nhất là Tư vấn pháp lý và Khôi phục phiên hội thoại nhằm minh họa luồng truyền thông điệp giữa các đối tượng.

% Sơ đồ trình tự mô tả cách thức hệ thống tiếp nhận một câu hỏi và trải qua các bước phân tích tác nhân để đưa ra câu trả lời.

% \begin{figure}[H]
%     \centering
%     \includegraphics[width=1\textwidth]{Hinhve/biểu đồ trình tự use case tư vấn pháp lý.png}
%     \caption{Biểu đồ trình tự luồng tư vấn pháp lý}
%     \label{fig:seq_consultation}
% \end{figure}

% Luồng thông điệp được thực hiện qua các bước: Người dùng gửi tin nhắn, sau đó Router Agent chọn retriever phù hợp với ý định câu hỏi của người dùng. Retriever truy xuất dữ liệu từ Neo4j, trả lại cho router agent. Router agent trả lại cho hàm main, tại đây main kết hợp context dữ liệu pháp lý truy vấn được với prompt rồi gửi cho LLM tổng hợp đáp án cuối cùng \ref{fig:seq_consultation}.

% \subsubsection{Use Case: Khôi phục phiên hội thoại (Resume Session)}
% Sơ đồ trình tự thể hiện sự tương tác giữa lớp trình diễn và lớp hạ tầng để đảm bảo tính liên tục của dữ liệu khi người dùng truy cập lại các hội thoại cũ.

% \begin{figure}[H]
%     \centering
%     % \includegraphics[width=0.85\textwidth]{figures/seq_resume.png}
%     \caption{Biểu đồ trình tự khôi phục phiên hội thoại}
%     \label{fig:seq_resume}
% \end{figure}

% Quy trình khôi phục bao gồm: Người dùng chọn phiên cũ $\rightarrow$ Hook on\_chat\_resume() kích hoạt $\rightarrow$ Adapter truy xuất lịch sử từ PostgreSQL $\rightarrow$ Presentation tái cấu trúc session\_history và hiển thị lại nội dung cho người dùng.

\subsection{Thiết kế chi tiết luồng Agentic Workflow}
\label{subsection:agentic-workflow}

Để giải quyết các yêu cầu phức tạp về pháp lý và đảm bảo tính chính xác của dữ liệu truy xuất, hệ thống được thiết kế theo mô hình Agentic GraphRAG, kết hợp Agentic RAG với Đồ thị tri thức (Knowledge Graph). Đây là một quy trình điều phối các tác nhân thông minh (Agents) chuyên biệt nhằm phân tích ý định người dùng, lựa chọn retriever agent phù hợp và tổng hợp câu trả lời tư vấn dựa trên căn cứ pháp lý lấy được từ Đồ thị tri thức.

\subsubsection{Sơ đồ luồng tổng quan}

Luồng hoạt động của hệ thống được chia thành 5 bước chính, được minh họa trong Hình \ref{fig:overview_flow}.

\begin{figure}[H]
    \centering
    \includegraphics[width=\textwidth]{Hinhve/Luồng Agentic GraphRAG.png}
    \caption{Sơ đồ luồng hoạt động Agentic GraphRAG}
    \label{fig:overview_flow}
\end{figure}

Cụ thể, Agentic Workflow gồm 5 bước như sau:
(i) \textbf{Bước 1 (Tiếp nhận câu hỏi và định tuyến)}: Chat Orchestrator tiếp nhận câu hỏi của người dùng cùng với lịch sử hội thoại của phiên chat hiện tại, sau đó kích hoạt Router Agent để phân tách ý hỏi và đề xuất danh sách các retriever agent cần gọi để trả lời câu hỏi. Danh sách các retriever agent được Router Agent đề xuất sẽ đi qua một bước kiểm tra để loại bỏ các tác nhân không tồn tại trong danh sách các retriever được đăng ký và loại bỏ các lời gọi trùng lặp tới cùng một retriever agent.
(ii) \textbf{Bước 2 (Phân tích ý định và trích xuất tham số)}: Mỗi Retriever Agent thực hiện hai nhiệm vụ: phân loại mẫu truy vấn Cypher (Cypher template) phù hợp với dạng câu hỏi và trích xuất các tham số cần thiết để điền vào mẫu truy vấn đã chọn, bao gồm cả mốc thời gian sự kiện pháp lý xuất hiện trong câu hỏi (nếu có).
(iii) \textbf{Bước 3 (Truy xuất căn cứ pháp lý)}: Các Retrieval Agent thực thi truy vấn Cypher với mẫu và tham số đã trích xuất để lấy về các căn cứ pháp lý cùng với các quan hệ liên quan như sửa đổi, hướng dẫn, thay thế hoặc mâu thuẫn.
(iv) \textbf{Bước 4 (Chuẩn hóa và hợp nhất ngữ cảnh)}: Chat Orchestrator gom các kết quả trả về từ nhiều Retriever Agents, loại bỏ các căn cứ pháp lý trùng lặp theo định danh và thời hạn hiệu lực, đồng thời sắp xếp ngữ cảnh thành các khối thông tin có cấu trúc gồm cảnh báo, hiệu lực văn bản, căn cứ chính, căn cứ hướng dẫn, căn cứ bổ trợ, căn cứ thay thế, căn cứ sắp có hiệu lực, căn cứ mâu thuẫn và các đường dẫn tham chiếu tới văn bản pháp luật. Khi ngữ cảnh pháp lý không có căn cứ phù hợp, hệ thống sẽ phản hồi không tìm thấy căn cứ và đề nghị người dùng đặt lại câu hỏi hoặc bổ sung thông tin.
(v) \textbf{Bước 5 (Sinh câu trả lời và hiển thị luồng xử lý)}: Luồng điều phối chat ghép ngữ cảnh đã chuẩn hóa với prompt chuyên biệt về lập luận pháp lý và lịch sử hội thoại, sau đó gọi API LLM để sinh câu trả lời. Câu trả lời cuối cùng có trích dẫn căn cứ pháp lý chi tiết tới cấp Điều, Khoản, Điểm; kèm theo các cảnh báo về hiệu lực văn bản và mâu thuẫn pháp lý (nếu có); đồng thời được hiển thị cùng với chi tiết luồng xử lý của các tác nhân và liên kết trực quan hóa phần đồ thị tri thức đã được truy xuất.

\subsubsection{Đặc tả chi tiết}

Trước tiên là đặc tả Chat Orchestrator tại Bảng \ref{table: chat_orchestrator}. Chat Orchestrator là thành phần điều phối luồng agentic workflow, đóng vai trò trung tâm kết nối các tác nhân, quản lý trạng thái phiên hội thoại và sinh câu trả lời cuối cùng dựa trên ngữ cảnh đã hợp nhất.

\begin{table}[H]
\centering
\begin{tabularx}{\textwidth}{|>{\raggedright\arraybackslash}p{3.8cm}|>{\raggedright\arraybackslash}X|}
\hline
\textbf{Thành phần} & \textbf{Mô tả chi tiết} \\ \hline
\textbf{Chức năng} & Khởi tạo và duy trì phiên hội thoại, kích hoạt Router Agent, thu thập kết quả truy xuất từ các Retrieval Agent, hợp nhất và chuẩn hóa ngữ cảnh pháp lý, gọi API LLM để sinh câu trả lời, đồng thời ghi lại vết thực thi của các tác nhân để hỗ trợ minh bạch hóa luồng xử lý cho người dùng. \\ \hline
\textbf{Đầu vào} & câu hỏi của người dùng; lịch sử hội thoại của phiên chat hiện tại; bộ nhớ ngữ cảnh pháp lý đã được các Retrieval Agent truy xuất ở các lượt hỏi đáp trước trong cùng phiên; phần ngữ cảnh pháp lý được tổng hợp từ kết quả của các Retrieval Agent trong lượt hỏi đáp hiện tại. \\ \hline
\textbf{Prompt} & Quy định chi tiết các quy tắc lập luận pháp lý và cách trích dẫn căn cứ pháp lý. \\ \hline
\textbf{Quy trình xử lý} & (i) Tiếp nhận câu hỏi (ii) Dựng lại lịch sử hội thoại (iii) Gọi Router Agent để chọn các retriever agent (iv) Thực thi song song các Retrieval Agent đã được chọn (v) Cập nhật bộ nhớ ngữ cảnh truy xuất của phiên (vi) Hợp nhất và loại bỏ căn cứ pháp lý trùng lặp theo định danh và thời hạn hiệu lực (vii) Dựng lịch sử hội thoại cho lượt sinh câu trả lời (viii) Gọi API LLM với prompt để sinh câu trả lời cuối cùng tới người dùng. \\ \hline
% \textbf{Quản lý trạng thái} & Khởi tạo trạng thái phiên hội thoại khi người dùng bắt đầu một đoạn chat mới; khôi phục lại lịch sử hội thoại và bộ nhớ ngữ cảnh khi người dùng mở lại phiên chat cũ; lưu vết toàn bộ luồng xử lý của các tác nhân (đầu vào, đầu ra) vào cơ sở dữ liệu lưu trữ phiên để phục vụ truy vết, kiểm thử và khôi phục hội thoại. \\ \hline
\textbf{Đầu ra} & Câu trả lời cho người dùng, link visualize đồ thị tri thức và các cảnh báo (nếu có). \\ \hline
\end{tabularx}
\caption{Đặc tả Chat Orchestrator}
\label{table: chat_orchestrator}
\end{table}

Tiếp theo là đặc tả Router Agent tại Bảng \ref{table: router_agent}. Tác nhân này chịu trách nhiệm phân tách ý hỏi của người dùng và lựa chọn các retriever agent phù hợp để giải quyết câu hỏi.

\begin{table}[H]
\centering
\begin{tabularx}{\textwidth}{|>{\raggedright\arraybackslash}p{3.8cm}|>{\raggedright\arraybackslash}X|}
\hline
\textbf{Thành phần} & \textbf{Mô tả chi tiết} \\ \hline
\textbf{Chức năng} & Phân loại ý định người dùng để gọi đúng và đủ các retriever agents cần thiết cho việc trả lời câu hỏi, viết lại câu hỏi của người người dùng dựa trên lịch sử hội thoại cho từng retriever \\ \hline
\textbf{Đầu vào} & Lịch sử hội thoại; câu truy vấn của người dùng; danh sách các retriever agent cùng mô tả chức năng và schema tham số tương ứng của từng retriever. \\ \hline
\textbf{Prompt} & Hướng dẫn LLM phân tách các câu hỏi đa ý của người dùng và xác định đầy đủ các bối cảnh pháp lý liên quan để gọi đủ retriever; tránh gọi trùng lặp cùng một retriever nhiều lần; viết lại câu hỏi từ người dùng dựa trên lịch sử hội thoại cho từng retriever.\\ \hline
\textbf{Quy trình xử lý} & (i) Nhận prompt cùng câu truy vấn, lịch sử hội thoại và danh sách retriever (ii) Gọi API LLM theo cơ chế tool calling (iii) Nhận về danh sách các retriever cần kích hoạt cùng với câu hỏi đã được viết lại cho mỗi lời gọi retriever. \\ \hline
\textbf{Quality gate} & Loại bỏ các retriever không tồn tại trong danh sách các retriever; loại bỏ các lời gọi trùng lặp tới cùng một retriever. Khi LLM không chọn được retriever phù hợp, hệ thống sử dụng retriever mặc định (fallback). \\ \hline
\textbf{Đầu ra} & Một mảng các lời gọi retriever, mỗi lời gọi gồm tên retriever và câu hỏi của người dùng được viết lại cho retriever đó. \\ \hline
\end{tabularx}
\caption{Đặc tả Router Agent}
\label{table: router_agent}
\end{table}

Cuối cùng là đặc tả retriever agent tại Bảng \ref{table: retriever_agent}. Tác nhân này có nhiệm vụ truy xuất Đồ thị tri thức để lấy về ngữ cảnh pháp lý theo từng dạng câu hỏi chuyên biệt phục vụ cho việc sinh câu trả lời.

\begin{table}[H]
\centering
\begin{tabularx}{\textwidth}{|>{\raggedright\arraybackslash}p{3.8cm}|>{\raggedright\arraybackslash}X|}
\hline
\textbf{Thành phần} & \textbf{Mô tả chi tiết} \\ \hline
\textbf{Chức năng} & Phân tích câu hỏi để lựa chọn (các) mẫu truy vấn Cypher (Cypher template) phù hợp, sau đó trích xuất các tham số rồi điền vào câu truy vấn Cypher để lấy về các căn cứ pháp lý. \\ \hline
\textbf{Đầu vào} & Câu truy vấn độc lập đã được viết lại bởi router agent. \\ \hline
\textbf{Prompt} & Yêu cầu LLM phân loại mẫu truy vấn Cypher cần dùng để truy xuất ngữ cảnh pháp lý dựa trên ý định pháp lý trong câu hỏi của người dùng và mô tả chức năng của từng mẫu truy vấn; trích xuất các tham số cần thiết và mốc thời gian sự kiện pháp lý xuất hiện trong câu hỏi. \\ \hline
\textbf{Quy trình xử lý} & (i) Phân loại mẫu truy vấn Cypher cần dùng (ii) Trích xuất các tham số truy vấn và mốc thời gian sự kiện pháp lý (ii) Thực thi truy vấn Cypher vào Đồ thị tri thức với mẫu và tham số đã trích xuất (iii) Chuẩn hóa ngữ cảnh pháp lý truy xuất, trả về cho Chat Orchestrator. \\ \hline
\textbf{Quality gate} & Khi LLM không chọn được mẫu truy vấn Cypher phù hợp, hệ thống sử dụng mẫu truy vấn Cypher mặc định (fallback). \\ \hline
\textbf{Đầu ra} & Danh sách các căn cứ pháp lý kèm các thuộc tính như nội dung điều luật, thời hạn hiệu lực, vai trò (căn cứ chính, căn cứ hướng dẫn, căn cứ bổ trợ, căn cứ thay thế, căn cứ sắp có hiệu lực, căn cứ mâu thuẫn). \\ \hline
\end{tabularx}
\caption{Đặc tả Retriever Agent}
\label{table: retriever_agent}
\end{table}

\subsection{Thiết kế cơ sở dữ liệu}
Ứng dụng sử dụng Neo4j để lưu trữ đồ thị tri thức Luật Hôn nhân Gia đình năm 2014 và hệ quản trị cơ sở dữ liệu PostgreSQL để lưu trữ user, lịch sử cuộc trò chuyện (Bao gồm các tin nhắn của người dùng, câu trả lời của chatbot và input/output của các agent trong quá trình trả lời)

Để thực hiện tính năng khôi phục phiên hội thoại và khả năng truy vết chi tiết luồng hoạt động của hệ thống Agentic GraphRAG, đồ án sử dụng hệ quản trị cơ sở dữ liệu quan hệ PostgreSQL để lưu trữ lịch sử hội thoại và thông tin người dùng. Việc thiết kế cơ sở dữ liệu tập trung vào khả năng quản lý phiên và lưu trữ từng bước (input/output) xử lý trung gian của các agent trong hệ thống.

Hình \ref{fig:er_diagram} thể hiện sơ đồ thực thể liên kết của hệ thống.

\begin{figure}[H]
    \centering
    \includegraphics[width=\textwidth]{Hinhve/ERdiagram.png}
    \caption{Sơ đồ thực thể liên kết cho chức năng lưu trữ hội thoại và truy vết luồng Agentic GraphRAG}
    \label{fig:er_diagram}
\end{figure}

Sơ đồ gồm sáu thực thể với các mối quan hệ như sau. Quan hệ Users -- Threads (một--nhiều): mỗi người dùng có thể tạo nhiều phiên hội thoại, mỗi phiên thuộc một user (định danh qua email hoặc OAuth). Quan hệ Users -- Projects (một--nhiều): mỗi user quản lý nhiều dự án nhóm hội thoại. Quan hệ Threads -- Steps (một--nhiều): mỗi phiên chứa chuỗi bước xử lý gồm tin nhắn user/assistant và vết thực thi agent (Router, Retriever). Quan hệ Steps -- Feedbacks (một--không hoặc một). Quan hệ Projects -- Project\_threads -- Threads: bảng liên kết gán tối đa một thread vào một project.

Bảng \ref{table:db_entities} liệt kê thực thể; các bảng tiếp theo mô tả chi tiết từng thực thể.

\begin{table}[H]
\centering
\begin{tabular}{|l|p{10cm}|}
\hline
\textbf{Tên thực thể} & \textbf{Ý nghĩa} \\ \hline
Users & Lưu trữ thông tin định danh và metadata tài khoản người dùng. \\ \hline
Threads & Quản lý các phiên hội thoại (sessions). \\ \hline
Steps & Lưu trữ tin nhắn, vết agent và metadata phục vụ khôi phục phiên. \\ \hline
Feedbacks & Lưu trữ đánh giá cho câu trả lời assistant. \\ \hline
Projects & Nhóm phiên hội thoại theo dự án do người dùng tạo. \\ \hline
Project\_threads & Liên kết thread với project (mỗi thread tối đa một project). \\ \hline
\end{tabular}
\caption{Danh sách các thực thể trong cơ sở dữ liệu}
\label{table:db_entities}
\end{table}

\begin{table}[H]
\centering
\begin{tabular}{|l|l|p{8cm}|}
\hline
\textbf{Tên trường} & \textbf{Kiểu dữ liệu} & \textbf{Mô tả} \\ \hline
id & UUID & Khóa chính, định danh duy nhất cho mỗi người dùng. \\ \hline
identifier & TEXT & Email hoặc thông tin đăng nhập từ nhà cung cấp OAuth (Google/Github). \\ \hline
metadata & JSONB & Lưu trữ thông tin bổ sung như tên hiển thị, ảnh đại diện dưới dạng JSON. \\ \hline
createdAt & TEXT & Thời điểm tạo tài khoản người dùng. \\ \hline
\end{tabular}
\caption{Thiết kế chi tiết bảng Users}
\label{table:db_users}
\end{table}

\begin{table}[H]
\centering
\begin{tabular}{|l|l|p{8cm}|}
\hline
\textbf{Tên trường} & \textbf{Kiểu dữ liệu} & \textbf{Mô tả} \\ \hline
id & UUID & Khóa chính, định danh duy nhất cho mỗi phiên hội thoại. \\ \hline
userId & UUID & Khóa ngoại, liên kết với thực thể Users để xác định chủ sở hữu phiên. \\ \hline
name & TEXT & Tên của phiên hội thoại (thường lấy theo nội dung câu hỏi đầu tiên). \\ \hline
createdAt & TEXT & Thời điểm khởi tạo phiên hội thoại. \\ \hline
\end{tabular}
\caption{Thiết kế chi tiết bảng Threads}
\label{table:db_threads}
\end{table}

\begin{table}[H]
\centering
\begin{tabular}{|l|l|p{8cm}|}
\hline
\textbf{Tên trường} & \textbf{Kiểu dữ liệu} & \textbf{Mô tả} \\ \hline
id & UUID & Khóa chính, định danh duy nhất cho từng bước xử lý. \\ \hline
threadId & UUID & Khóa ngoại liên kết với Threads. \\ \hline
parentId & UUID & ID step cha (cây step lồng nhau, ví dụ Router -- Retriever). \\ \hline
name & TEXT & Tên bước hoặc tác nhân (Router Agent, Retriever, Tổng hợp đáp án, \ldots). \\ \hline
type & TEXT & Loại bước: user\_message, assistant\_message, tool, llm, run. \\ \hline
input & TEXT & Đầu vào: câu hỏi, tham số retriever hoặc context gửi LLM. \\ \hline
output & TEXT & Đầu ra: tin nhắn hiển thị hoặc ngữ cảnh trích xuất. \\ \hline
metadata & JSONB & Lưu thông tin bổ sung theo từng bước: căn cứ pháp lý đã truy xuất, phiên bản đồ thị tri thức, cảnh báo hiệu lực pháp lý, các thông tin để xử lý lịch sử hội thoại \\ \hline
createdAt & TEXT & Thời điểm tạo step; dùng sắp xếp và khôi phục phiên. \\ \hline
\end{tabular}
\caption{Thiết kế chi tiết bảng Steps}
\label{table:db_steps}
\end{table}

\begin{table}[H]
\centering
\begin{tabular}{|l|l|p{8cm}|}
\hline
\textbf{Tên trường} & \textbf{Kiểu dữ liệu} & \textbf{Mô tả} \\ \hline
id & UUID & Khóa chính, định danh duy nhất cho mỗi feedback. \\ \hline
threadId & UUID & Khóa ngoại liên kết Threads. \\ \hline
forId & UUID & Khóa ngoại liên kết Steps. \\ \hline
value & INT & Nhận 2 giá trị 0 hoặc 1, 0 nếu Dislike, 1 nếu Like \\ \hline
comment & TEXT & comment chi tiết của người dùng về câu trả lời \\ \hline
\end{tabular}
\caption{Thiết kế chi tiết bảng Feedbacks}
\label{table:db_feedbacks}
\end{table}

\begin{table}[H]
\centering
\begin{tabular}{|l|l|p{8cm}|}
\hline
\textbf{Tên trường} & \textbf{Kiểu dữ liệu} & \textbf{Mô tả} \\ \hline
id & UUID & Khóa chính, định danh duy nhất cho mỗi project. \\ \hline
name & TEXT & Tên dự án hiển thị trên sidebar. \\ \hline
userId & UUID & Khóa ngoại liên kết Users (chủ sở hữu project). \\ \hline
createdAt & TEXT & Thời điểm tạo project. \\ \hline
updatedAt & TEXT & Thời điểm cập nhật gần nhất (tính cả hoạt động thread con). \\ \hline
\end{tabular}
\caption{Thiết kế chi tiết bảng Projects}
\label{table:db_projects}
\end{table}

\begin{table}[H]
\centering
\begin{tabular}{|l|l|p{8cm}|}
\hline
\textbf{Tên trường} & \textbf{Kiểu dữ liệu} & \textbf{Mô tả} \\ \hline
threadId & UUID & Khóa chính, khóa ngoại liên kết Threads. \\ \hline
projectId & UUID & Khóa ngoại liên kết Projects. \\ \hline
createdAt & TEXT & Thời điểm gán thread vào project. \\ \hline
\end{tabular}
\caption{Thiết kế chi tiết bảng Project\_threads}
\label{table:db_project_threads}
\end{table}

Về phần thiết kế đồ thị tri thức (Knowledge Graph) cho Luật hôn nhân gia đình năm 2014 sẽ được trình bày chi tiết ở chương \ref{chapter:5}

\section{Xây dựng ứng dụng}
\subsection{Thư viện và công cụ sử dụng}
Bảng \ref{table:my_label} liệt kê danh sách các thư viện và công cụ sử dụng để phát triển ứng dụng
\begin{table}[H]
\centering{}
    \begin{tabular}{|p{5cm}|p{3.5cm}|p{6.5cm}|}
        \hline
        \textbf{Mục đích} & \textbf{Công cụ} & \textbf{Địa chỉ URL} \\ \hline
        Ngôn ngữ lập trình & Python 3.13 & https://www.python.org/ \\ \hline
        Giao diện và điều phối & Chainlit 2.11.1 & https://chainlit.io/ \\ \hline
        Vẽ đồ thị tri thức & vis-network 9.1.9 & https://visjs.github.io/vis-network/ \\ \hline
        Framework LLM và agent & LangChain 1.2 & https://www.langchain.com/ \\ \hline
        Tích hợp API LLM & LangChain-OpenAI, LangChain-Google-GenAI & https://python.langchain.com/ \\ \hline
        Cơ sở dữ liệu đồ thị & Neo4j & https://neo4j.com/ \\ \hline
        Cơ sở dữ liệu quan hệ & PostgreSQL 15 & https://www.postgresql.org/ \\ \hline
        Trích xuất dữ liệu PDF & PyMuPDF 1.27 & https://pymupdf.readthedocs.io/ \\ \hline
        Ràng buộc dữ liệu & Pydantic 2.13 & https://docs.pydantic.dev/ \\ \hline
        ORM cho cơ sở dữ liệu & SQLAlchemy 2.0 & https://www.sqlalchemy.org/ \\ \hline
        Kết nối PostgreSQL & psycopg2-binary, asyncpg & https://www.postgresql.org/ \\ \hline
        Quản lý biến môi trường & python-dotenv & https://pypi.org/project/python-dotenv/ \\ \hline
        IDE lập trình & Visual Studio Code & https://code.visualstudio.com/ \\ \hline
        Đánh giá hệ thống & RAGAS & https://docs.ragas.io/ \\ \hline
        Triển khai cơ sở dữ liệu & Docker & https://www.docker.com/ \\ \hline
    \end{tabular}
    \caption{Danh sách thư viện và công cụ sử dụng}
    \label{table:my_label}
\end{table}
\subsection{Kết quả đạt được}
\label{subsection:4.7}
Sản phẩm chính của đồ án là ứng dụng chatbot tư vấn pháp luật chuyên sâu về lĩnh vực Hôn nhân và Gia đình tại Việt Nam được đóng gói dưới dạng một ứng dụng web hoàn chỉnh. Ứng dụng chatbot có khả năng giải đáp các câu hỏi lý thuyết và tình huống thực tế liên quan đến Luật Hôn nhân và gia đình 2014. Chatbot cho phép người dùng quan sát trực tiếp luồng hoạt động cũng như dữ liệu đầu vào và đầu ra của từng tác nhân AI (agent) tham gia vào quy trình truy xuất. Đồng thời, ĐATN cũng đã xây dựng được chức năng trực quan hóa phần đồ thị tri thức để trả lời cho câu hỏi của người dùng, và các cảnh báo về hiệu lực, mâu thuẫn (nếu có).

Dưới đây, bảng \ref{tab:summary_stats} chứa các thông tin thống kê của hệ thống:

\begin{table}[H]
\centering
\begin{tabular}{|l|l|}
\hline
\textbf{Thông tin} & \textbf{Giá trị} \\ \hline
Tổng số dòng mã nguồn & 50.818 dòng \\ \hline
Số file mã nguồn & 288 file \\ \hline
Dung lượng mã nguồn & khoảng 2,0 MB \\ \hline
Số retriever agent & 20 retriever \\ \hline
Số lượng câu hỏi trong tập testcase & 600 câu \\ \hline
Số lượng các node trong đồ thị tri thức & 14.971 node \\ \hline
Số lượng các relationship trong đồ thị tri thức & 16.867 relationship \\ \hline
\end{tabular}
\caption{Thông tin thống kê về ứng dụng}
\label{tab:summary_stats}
\end{table}

\subsection{Minh họa các chức năng chính}
\label{subsection:4.8}

Phần này trình bày các giao diện tiêu biểu minh họa cho những chức năng cốt lõi và đặc trưng nhất của hệ thống chatbot tư vấn pháp luật. Các tính năng được lựa chọn tập trung vào việc thể hiện tính minh bạch của kiến trúc Agentic GraphRAG và khả năng cung cấp căn cứ pháp lý chính xác cho người dùng.

\subsubsection{Chức năng theo dõi luồng hoạt động của tác nhân (Agent Trace)}

Một trong những tính năng quan trọng nhất của ứng dụng là khả năng hiển thị chi tiết quy trình xử lý nội bộ của các tác nhân thông minh. Chức năng này cho phép người dùng quan sát cách thức hệ thống bóc tách yêu cầu, lựa chọn công cụ và truy xuất dữ liệu từ đồ thị tri thức, giúp tăng cường độ tin cậy và tính xác thực cho các câu trả lời pháp lý.

\begin{figure}[H]
    \centering
    \includegraphics[width=1\textwidth]{Hinhve/ảnh màn hình luồng agent.png}
    \caption{Giao diện hiển thị luồng xử lý chi tiết và kết quả truy xuất của các Agent}
    \label{fig:agent_trace_ui}
\end{figure}

Giao diện minh họa tại Hình \ref{fig:agent_trace_ui} thể hiện sự phân tách rõ ràng giữa các bước tư duy của máy và kết quả tư vấn cuối cùng. Trong phân đoạn "Luồng truy xuất ngữ cảnh", hệ thống ghi lại vết thực thi của tác nhân \textit{Router Agent} khi định tuyến đến bộ retriever agent dữ liệu phù hợp. Tiếp đến, tại phân đoạn "Tổng hợp đáp án", người dùng có thể kiểm tra trực tiếp ngữ cảnh pháp lý (context) được trích xuất từ cơ sở dữ liệu Neo4j để đối chiếu tính chính xác của thông tin nguồn. Sau cùng là câu trả lời từ chatbot.

\subsubsection{Chức năng trực quan hóa đồ thị tri thức}

\begin{figure}[H]
    \centering
    \includegraphics[width=1\textwidth]{Hinhve/ảnh màn hình visualize đồ thị tri thức.png}
    \caption{Giao diện hiển thị đồ thị tri thức}
    \label{fig:visualize_graph_ui}
\end{figure}

Giao diện minh họa tại Hình \ref{fig:visualize_graph_ui} thể hiện sự trực quan hóa đồ thị tri thức. Người dùng có thể quan sát trực tiếp các node và relationship trong đồ thị tri thức, tương tác với các node trên đồ thị tri thức bằng các thao tác thu, phóng, kéo thả, tìm kiếm. Người dùng cũng có thể xem template được chọn, và các tham số đã được trích xuất để điền vào cypher query. Người dùng cũng có thể xem danh sách các node và quan hệ giữa chúng ở dạng bảng và JSON.

\subsubsection{Chức năng cảnh báo hiệu lực, mâu thuẫn pháp lý}

\begin{figure}[H]
    \centering
    \includegraphics[width=1\textwidth]{Hinhve/ảnh màn hình cảnh báo hiệu lực.png}
    \caption{Giao diện modal cảnh báo hiệu lực}
    \label{fig:legal_warning_ui}
\end{figure}

\begin{figure}[H]
    \centering
    \includegraphics[width=1\textwidth]{Hinhve/ảnh màn hình cảnh báo mâu thuẫn.png}
    \caption{Giao diện modal cảnh báo mâu thuẫn}
    \label{fig:legal_warning_ui_2}
\end{figure}

Giao diện modal cảnh báo hiệu lực và mâu thuẫn pháp lý được hiển thị khi người dùng tương tác với đồ thị tri thức tại Hình \ref{fig:legal_warning_ui} và Hình \ref{fig:legal_warning_ui_2}. Người dùng có thể xem các căn cứ pháp lý đã được truy xuất, và các cảnh báo về hiệu lực, mâu thuẫn (nếu có).

\section{Kiểm thử}
\label{section:4.9}

Nội dung chương này tập trung vào việc xác định các kỹ thuật kiểm thử, thiết lập bộ tiêu chí đánh giá định lượng và thực hiện các kịch bản kiểm thử dựa trên các nhóm câu hỏi pháp lý chuyên biệt.

\subsection{Kỹ thuật kiểm thử và tiêu chí đánh giá}

Đồ án kết hợp hai phương pháp kiểm thử chính để đánh giá toàn diện hệ thống. Thứ nhất là kỹ thuật kiểm thử hộp đen (Black-box testing) được tiến hành một cách thủ công nhằm đánh giá khả năng phản hồi ngôn ngữ tự nhiên và tính đúng đắn của logic tư vấn đối với người dùng cuối. Thứ hai là kỹ thuật đánh giá tự động dựa trên thư viện RAGAS nhằm đo lường chất lượng của quy trình truy xuất và sinh câu trả lời dựa trên ngữ cảnh pháp lý. Việc đánh giá được thực hiện trên bộ testcase gồm 600 câu hỏi/câu trả lời, chi tiết về cách xây dựng bộ testcase sẽ được trình bày ở chương \ref{chapter:5}

Bảng \ref{table:testing_metrics} liệt kê các tiêu chí vả chỉ số đánh giá hệ thống.
\begin{table}[H]
\centering
\caption{Tiêu chí và chỉ số đánh giá hiệu năng hệ thống}
\label{table:testing_metrics}
\begin{tabular}{|l|l|p{4cm}|}
\hline
\textbf{Nhóm tiêu chí} & \textbf{Chỉ số đánh giá (Metrics)} & \textbf{Mục tiêu đạt được} \\ \hline
\multirow{3}{*}{Metric RAGAS} & Faithfulness & Đạt mức $>$ 85\% \\ \cline{2-3} 
& Context Recall & Đạt mức $>$ 85\% \\ \hline
\multirow{2}{*}{Căn cứ pháp lý} & Legal Validity Accuracy & Đạt mức $>$ 90\% \\ \cline{2-3} 
 & Legal Citation Accuracy & Đạt mức $>$ 85\% \\ \hline
\multirow{2}{*}{Xử lý xung đột} & Legal Conflict Handling Accuracy & Đạt mức $>$ 80\% \\ \cline{2-3} 
 & Context Conflict Detection Rate & Đạt mức $>$ 80\% \\ \hline
\multirow{2}{*}{Phân loại viện dẫn} & Citation Prioritization Accuracy & Đạt mức $>$ 85\% \\ \cline{2-3} 
 & Invalid Citation Rate & Đạt mức dưới 10\%  \\ \hline
\end{tabular}
\end{table}

\subsection{Chi tiết cách tính từng chỉ số}
\label{subsection:4.chi_tiet_chi_so}

Trước tiên là nhóm chỉ số RAGAS. Nhóm chỉ số này sử dụng LLM để đánh giá, và được định nghĩa chính xác theo mô tả của thư viện RAGAS như sau:

Chỉ số Faithfulness (Độ trung thực) tập trung đo lường tỷ lệ các tuyên bố trong câu trả lời có thể được suy diễn trực tiếp và hoàn toàn từ ngữ cảnh trích xuất từ đồ thị tri thức.

Chỉ số Context Recall (Độ phủ ngữ cảnh) chịu trách nhiệm đo lường mức độ đồng thuận giữa các căn cứ pháp lý thực tế được hệ thống gọi ra so với dữ liệu chuẩn Ground Truth của tập testcase được xây dựng.

Tiếp theo là nhóm chỉ số Đánh giá Căn cứ Pháp lý. Nhóm chỉ số này tập trung kiểm soát tính chính xác của các căn cứ pháp lý được viện dẫn trong câu trả lời của chatbot.

Chỉ số Legal Citation Accuracy (LCA -- Độ chính xác trích dẫn pháp lý) đo lường độ chính xác mã hiệu các Điều, Khoản văn bản luật cần thiết để giải quyết tình huống câu hỏi. Trong đó, LCA là viết tắt của \textit{Legal Citation Accuracy}.
\begin{equation}
\label{eq:lca_metric}
\text{LCA} = \frac{\text{Số điều luật trích dẫn đúng}}{\text{Tổng số điều luật cần trích dẫn (Ground Truth)}}
\end{equation}

Chỉ số Legal Validity Accuracy (LVA -- Độ chính xác về hiệu lực pháp luật) thực hiện đo lường khả năng kiểm soát dòng thời gian để loại bỏ các điều luật đã hết hiệu lực thi hành hoặc không phù hợp tại thời điểm xảy ra sự kiện của câu hỏi. Trong đó, LVA là viết tắt của \textit{Legal Validity Accuracy}.
\begin{equation}
\label{eq:lva_metric}
\text{LVA} = \frac{\text{Tổng số căn cứ còn hiệu lực được hệ thống gọi ra}}{\text{Tổng số căn cứ hệ thống đã gọi ra (trong toàn bộ testcase)}}
\end{equation}

Tiếp theo là nhóm chỉ số đánh giá khả năng xử lý Xung đột Pháp luật. Nhóm chỉ số này đánh giá khả năng phát hiện và xử lý áp dụng các căn cứ theo đúng nguyên tắc thứ bậc hiệu lực pháp lý khi tình huống có sự chồng chéo, mâu thuẫn.

Chỉ số Context Conflict Detection Rate (CCDR -- Tỷ lệ phát hiện xung đột trong ngữ cảnh) đo lường khả năng trong việc nhận diện và đánh dấu (\textit{flag}) các điều luật chồng chéo hoặc mâu thuẫn. Trong đó, CCDR là viết tắt của \textit{Context Conflict Detection Rate}. Chỉ số này được tính như sau:
\begin{equation}
\label{eq:ccdr_metric}
\text{CCDR} = \frac{\text{Số câu phát hiện và đánh dấu (flag) có xung đột đúng}}{\text{Tổng số câu thực sự có xung đột trong tập testcase}}
\end{equation}

Chỉ số Legal Conflict Handling Accuracy (LCHA -- Độ chính xác xử lý xung đột pháp luật) đánh giá khả năng áp dụng các nguyên tắc ưu tiên thứ bậc hiệu lực pháp lý để đưa ra đáp án đúng cuối cùng. Trong đó, LCHA là viết tắt của \textit{Legal Conflict Handling Accuracy}. Công thức tính toán như sau:
\begin{equation}
\label{eq:lcha_metric}
\text{LCHA} = \frac{\text{Số câu hỏi chồng chéo, mâu thuẫn được xử lý đúng nguyên tắc ưu tiên}}{\text{Tổng số câu hỏi có chứa xung đột pháp lý trong testcase}}
\end{equation}

Tiếp theo là nhóm chỉ số Phân loại viện dẫn. Nhóm chỉ số này đo lường độ chính xác của hệ thống trong việc phân rã vai trò ngữ nghĩa của các tầng văn bản luật khi thực hiện trích xuất tri thức.

Chỉ số Citation Prioritization Accuracy (CPA -- Độ chính xác phân cấp viện dẫn) đóng vai trò đánh giá năng lực của hệ thống trong việc sắp xếp trật tự và gán nhãn chính xác vai trò văn bản bao gồm căn cứ chính, căn cứ bổ trợ, hay văn bản hướng dẫn, sửa đổi, thay thế. Trong đó, CPA là viết tắt của \textit{Citation Prioritization Accuracy}. Chỉ số được tính như sau: 
\begin{equation}
\label{eq:cpa_metric}
\text{CPA} = \frac{1}{3} \times \left( \begin{aligned}
    &\frac{\text{Số căn cứ chính đúng}}{\text{Tổng số căn cứ pháp lý chính}} \\
    + &\frac{\text{Số căn cứ bổ trợ đúng}}{\text{Tổng số căn cứ pháp lý bổ trợ}} \\
    + &\frac{\text{Số căn cứ hướng dẫn, thay thế đúng}}{\text{Tổng số căn cứ hướng dẫn, sửa đổi, bổ sung, thay thế}}
\end{aligned} \right)
\end{equation}

Chỉ số Invalid Citation Rate (ICR -- Tỷ lệ trích dẫn sai phạm vi) thực hiện đo lường tỷ lệ xuất hiện lỗi nghiêm trọng khi hệ thống gọi ra các điều khoản không liên quan đến câu hỏi của người dùng, đã hết hiệu lực, không phù hợp với thời điểm xảy ra sự kiện pháp lý trong câu hỏi của người dùng hoặc đã bị bãi bỏ hoàn toàn. Trong đó, ICR là viết tắt của \textit{Invalid Citation Rate}. Công thức tính toán như sau:
\begin{equation}
\label{eq:icr_metric}
\text{ICR} = \frac{\text{Số câu trả lời chứa ít nhất 1 căn cứ hết hiệu lực hoặc không liên quan tình huống}}{\text{Tổng số câu hỏi trong bộ testcase}}
\end{equation}

\subsection{Kịch bản kiểm thử theo dạng câu hỏi pháp lý}
\label{subsection:4.9.2}
% \subsubsection{Kịch bản 1: Nhóm câu hỏi về điều kiện kết hôn}
% \begin{table}[H]
% \centering
% \begin{tabular}{|l|p{10cm}|}
% \hline
% \textbf{Thành phần} & \textbf{Mô tả chi tiết kịch bản} \\ \hline
% Mục tiêu & Kiểm tra khả năng tư vấn về độ tuổi, sự tự nguyện và các trường hợp cấm kết hôn. \\ \hline
% Dạng câu hỏi & "Nam 19 tuổi và nữ 18 tuổi có đủ điều kiện đăng ký kết hôn không?" \\ \hline
% Kết quả mong đợi & Hệ thống xác định đúng Điều 8 Luật Hôn nhân gia đình 2014 và chỉ rõ điều kiện về độ tuổi. \\ \hline
% \end{tabular}
% \caption{Kịch bản kiểm thử nhóm điều kiện kết hôn}
% \end{table}
Bộ dữ liệu kiểm thử gồm 600 câu hỏi trong bộ testcase được thiết kế chia ra thành các dạng câu hỏi là 
(i) Điều kiện kết hôn;
(ii) Đăng ký kết hôn;
(iii) Kết hôn trái pháp luật;
(iv) Chung sống như vợ chồng;
(v) Ly hôn (Các quy định chung và quyền yêu cầu ly hôn);
(vi) Chia tài sản sau khi ly hôn;
(vii) Cha mẹ con sau ly hôn (Quyền và nghĩa vụ cha, mẹ với con sau khi ly hôn);
(viii) Chế độ tài sản của vợ chồng;
(ix) Xử phạt vi phạm (Chế tài và hình thức xử phạt vi phạm trong lĩnh vực hôn nhân);
(x) Cấp dưỡng (Nghĩa vụ cấp dưỡng giữa các thành viên trong gia đình);
(xi) Hôn nhân chấm dứt do vợ, chồng chết (Chấm dứt hôn nhân do vợ hoặc chồng chết hoặc bị tuyên bố là đã chết);
(xii) Hôn nhân có yếu tố nước ngoài;
(xiii) Các quy định chung, khái niệm pháp lý (Khái niệm pháp lý nền tảng và nguyên tắc của chế độ hôn nhân gia đình);
(xiv) Tài sản riêng của con (Quyền sở hữu và định đoạt tài sản riêng của con chưa thành niên);
(xv) Xác định cha, mẹ, con;

\subsection{Tổng kết kết quả kiểm thử}
\label{subsection:4.9.3}

Kết quả thực nghiệm được trình bày theo từng dạng câu hỏi pháp lý được nêu ở mục \ref{subsection:4.9.2}. Bảng \ref{table:experimental_results} trình bày kết quả chi tiết các chỉ số căn cứ pháp lý và phân loại viện dẫn theo từng chủ đề; Bảng \ref{table:ragas_results} trình bày kết quả hai chỉ số RAGAS (Faithfulness và Context Recall) trên cùng bộ 200 câu hỏi; sau đó, Bảng \ref{table:experimental_summary} tổng hợp các chỉ số đánh giá trung bình trên toàn bộ các câu hỏi thuộc bộ testcase.

\begin{table}[H]
\centering
\small
\begin{tabular}{|c|p{3cm}|c|c|c|c|c|c|c|}
\hline
\textbf{STT} & \textbf{Chủ đề pháp lý} & \textbf{Số câu} & \textbf{LT} & \textbf{TH} & \textbf{LCA} & \textbf{LVA} & \textbf{CPA} & \textbf{ICR} \\ \hline

1 & Điều kiện kết hôn & 18 & 14 & 4 & 0.8667 & 1.0 & 0.8889 & 0.0556 \\ \hline
2 & Đăng ký kết hôn & 12 & 9 & 3 & 0.8750 & 1.0 & 0.8750 & 0.0833 \\ \hline
3 & Kết hôn trái pháp luật & 10 & 8 & 2 & 0.8000 & 1.0 & 0.8500 & 0.3000 \\ \hline
4 & Chung sống như vợ chồng & 14 & 11 & 3 & 0.8214 & 1.0 & 0.8714 & 0.1429 \\ \hline
5 & Ly hôn & 22 & 17 & 5 & 0.8387 & 1.0 & 0.8773 & 0.1364 \\ \hline
6 & Chia tài sản sau ly hôn & 20 & 15 & 5 & 0.8250 & 1.0 & 0.8700 & 0.1500 \\ \hline
7 & Cha, mẹ, con sau ly hôn & 7 & 7 & 0 & 0.9057 & 1.0 & 0.8571 & 0.0 \\ \hline
8 & Chế độ tài sản của vợ chồng & 28 & 21 & 7 & 0.8182 & 1.0 & 0.8914 & 0.0 \\ \hline
9 & Xử phạt vi phạm & 8 & 6 & 2 & 0.8750 & 1.0 & 0.8875 & 0.1250 \\ \hline
10 & Cấp dưỡng & 6 & 6 & 0 & 0.8333 & 1.0 & 0.8333 & 0.0 \\ \hline
11 & Hôn nhân chấm dứt do vợ, chồng chết & 3 & 2 & 1 & 0.8333 & 1.0 & 0.89 & 0.0 \\ \hline
12 & Hôn nhân có yếu tố nước ngoài & 12 & 9 & 3 & 0.7917 & 1.0 & 0.8583 & 0.3333 \\ \hline
13 & Các quy định chung, khái niệm pháp lý & 16 & 13 & 3 & 0.8750 & 1.0 & 0.9125 & 0.0625 \\ \hline
14 & Tài sản riêng của con chưa thành niên & 10 & 8 & 2 & 0.8600 & 1.0 & 0.8800 & 0.1000 \\ \hline
15 & Xác định cha, mẹ, con & 14 & 11 & 3 & 0.7857 & 1.0 & 0.8643 & 0.3571 \\ \hline

\multicolumn{2}{|c|}{\textbf{Tổng cộng / Trung bình}} & \textbf{200} & \textbf{157} & \textbf{43} & \textbf{0.8370} & \textbf{1.0} & \textbf{0.878} & \textbf{0.125} \\ \hline
\end{tabular}
\caption{Kết quả kiểm thử thực nghiệm chi tiết theo từng chủ đề pháp lý}
\label{table:experimental_results}
\end{table}

\begin{table}[H]
\centering
\small
\begin{tabular}{|c|p{3cm}|c|c|c|c|c|}
\hline
\textbf{STT} & \textbf{Chủ đề pháp lý} & \textbf{Số câu} & \textbf{LT} & \textbf{TH} & \textbf{Faithfulness} & \textbf{Context Recall} \\ \hline

1 & Điều kiện kết hôn & 18 & 14 & 4 & 0.9189 & 0.9172 \\ \hline
2 & Đăng ký kết hôn & 12 & 9 & 3 & 0.9240 & 0.9100 \\ \hline
3 & Kết hôn trái pháp luật & 10 & 8 & 2 & 0.8200 & 0.7900 \\ \hline
4 & Chung sống như vợ chồng & 14 & 11 & 3 & 0.8714 & 0.8429 \\ \hline
5 & Ly hôn & 22 & 17 & 5 & 0.8818 & 0.8636 \\ \hline
6 & Chia tài sản sau ly hôn & 20 & 15 & 5 & 0.8750 & 0.8500 \\ \hline
7 & Cha, mẹ, con sau ly hôn & 7 & 7 & 0 & 0.9286 & 0.9286 \\ \hline
8 & Chế độ tài sản của vợ chồng & 28 & 21 & 7 & 0.8869 & 0.8736 \\ \hline
9 & Xử phạt vi phạm & 8 & 6 & 2 & 0.9125 & 0.9000 \\ \hline
10 & Cấp dưỡng & 6 & 6 & 0 & 0.9000 & 0.8833 \\ \hline
11 & Hôn nhân chấm dứt do vợ, chồng chết & 3 & 2 & 1 & 0.8833 & 0.8667 \\ \hline
12 & Hôn nhân có yếu tố nước ngoài & 12 & 9 & 3 & 0.8167 & 0.7917 \\ \hline
13 & Các quy định chung, khái niệm pháp lý & 16 & 13 & 3 & 0.9288 & 0.9213 \\ \hline
14 & Tài sản riêng của con chưa thành niên & 10 & 8 & 2 & 0.8900 & 0.8700 \\ \hline
15 & Xác định cha, mẹ, con & 14 & 11 & 3 & 0.8000 & 0.7786 \\ \hline

\multicolumn{2}{|c|}{\textbf{Tổng cộng / Trung bình}} & \textbf{200} & \textbf{157} & \textbf{43} & \textbf{0.8820} & \textbf{0.8650} \\ \hline
\end{tabular}
\caption{Kết quả kiểm thử các chỉ số RAGAS theo từng chủ đề pháp lý}
\label{table:ragas_results}
\end{table}

\begin{table}[H]
\centering
\begin{tabular}{|l|l|l|l|}
\hline
\textbf{Chỉ số đánh giá} & \textbf{Mục tiêu} & \textbf{Kết quả đạt được} & \textbf{Trạng thái} \\ \hline
Faithfulness & $>$ 85\% & 88.2\% & Đạt \\ \hline
Context Recall & $>$ 85\% & 86.5\% & Đạt \\ \hline
Legal Citation Accuracy & $>$ 85\% & 83.7\% & Không \\ \hline
Legal Validity Accuracy & $>$ 90\% & 100\% & Đạt \\ \hline
Citation Prioritization Accuracy & $>$ 85\% & 87.8\% & Đạt \\ \hline
Invalid Citation Rate & $<$ 10\% & 12.5\% & Không \\ \hline
Legal Conflict Handling Accuracy & $>$ 80\% & 82.4\% & Đạt \\ \hline
Context Conflict Detection Rate  & $>$ 80\% & 83.1\% & Đạt \\ \hline
\end{tabular}
\caption{Tổng hợp kết quả thực nghiệm và đối chiếu mục tiêu}
\label{table:experimental_summary}
\end{table}

Các nguyên nhân dẫn đến các trường hợp không đạt bao gồm: (i) sự thiếu hụt các mối quan hệ dẫn chiếu trong đồ thị tri thức, (ii) giới hạn về khả năng hiểu ngữ cảnh của mô hình ngôn ngữ lớn khi xử lý các câu hỏi có tình huống phức tạp, các câu hỏi nhiều ý hoặc (iii) việc phân loại các chủ đề pháp lý vẫn chưa đủ rõ ràng khiến router agent chọn sai retriever. Các hạn chế này sẽ là cơ sở quan trọng để đề xuất hướng cải thiện trong tương lai.

\section{Triển khai}
Hệ thống được đóng gói bằng Docker và triển khai trên nền tảng Render.com qua cơ chế Blueprint. 

Bảng \ref{table:render_infra} liệt kê cấu hình các dịch vụ trên Render được lựa chọn cho hệ thống.

\begin{table}[H]
\centering
\small
\setlength{\tabcolsep}{6pt} % Tăng nhẹ khoảng cách cột cho thoáng vì đã bớt cột
\begin{tabularx}{\textwidth}{|p{3cm}|X|p{3.5cm}|}
\hline
\textbf{Dịch vụ} & \textbf{Loại Render} & \textbf{Cấu hình} \\ \hline
app (Chainlit) & Web Service (Standard) & 2 GB RAM / 1 vCPU \\ \hline
neo4j & Private Service (Standard), Disk 5 GB & 2 GB RAM / 1 vCPU \\ \hline
db & Managed PostgreSQL (Starter) & 256 MB shared, 1 GB storage \\ \hline
\end{tabularx}
\caption{Cấu hình hạ tầng triển khai ứng dụng trên Render}
\label{table:render_infra}
\end{table}

Người dùng cũng có thể chạy ứng dụng trên máy bằng cách làm theo 2 bước sau đây:

Bước 1: Cài đặt Docker Desktop phiên bản 24 trở lên trên máy phát triển, sau đó clone mã nguồn của hệ thống. Người dùng đặt tệp backup đồ thị được cung cấp và đặt vào thư mục deploy/neo4j/snapshots/.

Bước 2: Sao chép tệp .env.example thành .env và điền OPENAI\_API\_KEY, GOOGLE\_API\_KEY, CHAINLIT\_AUTH\_SECRET (chuỗi ngẫu nhiên 32 ký tự). Khởi chạy ứng dụng bằng docker compose up -d -{}-build. Người dùng truy cập \url{http://localhost:8000} và bắt đầu sử dụng ứng dụng.

% \textbf{Giai đoạn 3 -- Đẩy mã nguồn lên GitHub và tạo Blueprint trên Render.} Khởi tạo một repository riêng tư trên GitHub (không công khai do snapshot có thể chứa dữ liệu nội bộ), commit toàn bộ mã nguồn rồi push lên nhánh main. Tại bảng điều khiển Render, đăng nhập và liên kết tài khoản GitHub, chọn \textit{New $\rightarrow$ Blueprint $\rightarrow$ repository vừa tạo}. Render tự động đọc tệp render.yaml ở thư mục gốc, xác định ba dịch vụ cần khởi tạo và yêu cầu xác nhận trước khi triển khai.

% \textbf{Giai đoạn 4 -- Cấu hình các biến môi trường bí mật.} Sau khi Blueprint được áp dụng, người dùng vào trang quản trị của dịch vụ app, mục \textit{Environment} và điền các biến môi trường bí mật theo Bảng \ref{table:render_secrets}. Các biến môi trường còn lại như NEO4J\_URI, DATABASE\_URL, NEO4J\_PASSWORD được Render tự động liên kết giữa các dịch vụ qua khai báo fromService và fromDatabase trong Blueprint, không cần thiết lập thủ công.

% \begin{table}[H]
% \centering
% \small
% \setlength{\tabcolsep}{5pt}
% \begin{tabularx}{0.92\textwidth}{|p{5.2cm}|X|}
% \hline
% \textbf{Biến môi trường} & \textbf{Cách điền} \\ \hline
% OPENAI\_API\_KEY & Dán khóa API lấy từ \url{https://platform.openai.com/api-keys} \\ \hline
% GOOGLE\_API\_KEY & Dán khóa API Gemini lấy từ Google AI Studio, dùng cho mô hình embedding \\ \hline
% CHAINLIT\_AUTH\_SECRET & Render tự sinh giá trị ngẫu nhiên, không cần chỉnh sửa \\ \hline
% OAUTH\_GOOGLE\_\newline CLIENT\_ID/SECRET & (Tùy chọn) Khai báo trên Google Cloud Console nếu cần bật đăng nhập OAuth \\ \hline
% \end{tabularx}
% \caption{Các biến môi trường bí mật cần điền thủ công trên Render}
% \label{table:render_secrets}
% \end{table}

% \textbf{Giai đoạn 5 -- Xác minh sau triển khai.} Render tiến hành build các image Docker và khởi chạy các dịch vụ trong khoảng năm đến mười phút. Người dùng theo dõi log của dịch vụ neo4j cho đến khi xuất hiện thông báo \textit{Database restored}, sau đó theo dõi log dịch vụ app đến khi xuất hiện thông báo \textit{Your app is available at}. Khi hoàn tất, hệ thống có thể truy cập công khai qua đường dẫn HTTPS https://[tên-app].onrender.com. Trong trường hợp cần sử dụng tên miền riêng, người dùng vào mục \textit{Settings $\rightarrow$ Custom Domains} của dịch vụ app, khai báo tên miền và cập nhật bản ghi CNAME tại nhà cung cấp DNS theo hướng dẫn của Render.

\end{document}
